\documentclass[11pt]{article}
\usepackage[utf8]{inputenc}
\usepackage[T1]{fontenc}
\usepackage{amsmath,amssymb,amsthm,mathtools}
\usepackage[margin=1in]{geometry}
\usepackage{xcolor,enumitem,hyperref,listings,booktabs}

\definecolor{navy}{RGB}{25,55,95}
\definecolor{lightblue}{RGB}{235,243,252}
\hypersetup{colorlinks=true,linkcolor=navy,urlcolor=blue!60!black}
\setlist{itemsep=3pt,topsep=4pt}
\setlength{\parindent}{0pt}
\setlength{\emergencystretch}{3em}
\newcommand{\R}{\mathbb{R}}
\newcommand{\C}{\mathbb{C}}
\newcommand{\norm}[1]{\left\lVert #1\right\rVert}
\newcommand{\T}{^{\mathsf T}}
\newcommand{\Hh}{^{\mathsf H}}
\newcommand{\tr}{\operatorname{tr}}
\newcommand{\diag}{\operatorname{diag}}
\newcommand{\Span}{\operatorname{span}}
\newcommand{\Null}{\operatorname{Null}}
\newcommand{\im}{\operatorname{im}}
\newcommand{\rank}{\operatorname{rank}}
\newcommand{\N}{\mathcal{N}}
\newtheorem{theorem}{Theorem}
\newtheorem{lemma}{Lemma}
\newtheorem{remark}{Remark}

\lstset{language=Python,basicstyle=\ttfamily\small,keywordstyle=\color{blue}\bfseries,
  commentstyle=\color{green!45!black}\itshape,stringstyle=\color{red!65!black},
  showstringspaces=false,numbers=left,numberstyle=\tiny\color{gray},
  breaklines=true,frame=single,framerule=0.5pt,framesep=4pt,xleftmargin=2em}

\title{\color{navy}\bfseries SLAI MF2601: Homework I}
\author{Name: \underline{\hspace{6cm}} \quad Student ID: \underline{\hspace{3cm}}}
\date{Due: September 30, 2026}

\begin{document}
\maketitle


\textbf{General instructions.} Unless a problem says otherwise, all matrices are real;
use MATLAB or Python for numerical work. Use \LaTeX{} to typeset the homework and
compile to PDF; you may modify the source code provided. Submit a single zip file
containing the PDF (with the theory and numerical-experiment results) together with
a MATLAB or Python code file, to the TA's email with subject line \texttt{NLA HW1}.

\section*{[Do not submit this part] Problem 0: Preliminary questions}

This is a study aid only. Do not submit answers to these questions.
You may use any resources, including AI tools, to answer these questions to
ensure a \textbf{basic} understanding of linear algebra.

\begin{quote}
What is the definition of \textbf{linear space} or vector space?\par
What is the definition of a \textbf{basis} of a linear space?\par
What is the definition of the \textbf{dimension} of a linear space?\par
What is the definition of \textbf{linear mapping}?\par
What is the definition of the \textbf{image and kernel} of a linear mapping?\par
What is the \textbf{relationship between matrix and linear mapping} under a fixed basis? \par
What happens to the matrix of a linear mapping if we change the basis?\par
What happens to their matrices if two linear mappings are composed?\par
What is the \textbf{rank} and \textbf{nullity} of a matrix?\par
Kernel -- nullity. Image -- rank. What are their relationships?\par
What is the \textbf{rank-nullity} theorem?\par
\par
\medskip
What is the definition of a linear \textbf{subspace} and the summation and intersection
between two subspaces?\par
How do you prove the dimension formula for the sum of two subspaces?
\[
\dim(S+T)=\dim S+\dim T-\dim(S\cap T)
\]
What is the definition of inner product?\par
What is the definition of orthogonal vectors?
\end{quote}


\section*{Problem 1}
\textbf{Page 16, Problem 2.6.} If $u$ and $v$ are $m$-vectors, the matrix $A = I + uv\Hh$ is known as a
\emph{rank-one perturbation of the identity}. Show that if $A$ is nonsingular, then its inverse
has the form $A^{-1} = I + \alpha uv\Hh$ for some scalar $\alpha$, and give an expression for
$\alpha$. For what $u$ and $v$ is $A$ singular? If it is singular, what is $\mathrm{null}(A)$?

\section*{Problem 2}
\textbf{Page 24, Problem 3.6.} Let $\|\cdot\|$ denote any norm on $\mathbb{C}^m$. The corresponding
\emph{dual norm} $\|\cdot\|'$ is defined by the formula
\[
\|x\|' = \sup_{\|y\| = 1} |y\Hh x|.
\]
\begin{enumerate}
\item[(a)] Prove that $\|\cdot\|'$ is a norm.
\item[(b)] Let $x, y \in \mathbb{C}^m$ with $\|x\| = \|y\| = 1$ be given. Show that there exists a
rank-one matrix $B = yz\Hh$ such that $Bx = y$ and $\|B\| = 1$, where $\|B\|$ is the matrix
norm of $B$ induced by the vector norm $\|\cdot\|$. You may use the following lemma,
without proof: given $x \in \mathbb{C}^m$, there exists a nonzero $z \in \mathbb{C}^m$ such
that $|z\Hh x| = \|z\|' \|x\|$.
\end{enumerate}

\section*{Problem 3}
\textbf{Page 31, Problem 4.4.} Two matrices $A, B \in \mathbb{C}^{m \times m}$ are \emph{unitarily
equivalent} if $A = QBQ\Hh$ for some unitary $Q \in \mathbb{C}^{m \times m}$. Is it true or
false that $A$ and $B$ are unitarily equivalent if and only if they have the same multiset
of singular values, counted with multiplicity? \textit{If yes, please give the proof; otherwise
provide an example.}

\section*{Problem 4}
\textbf{Page 37, Problem 5.4.} Suppose $A \in \mathbb{C}^{m \times m}$ has an SVD $A = U\Sigma V\Hh$.
Find an eigenvalue decomposition (5.1) of the $2m \times 2m$ Hermitian matrix
\[
\begin{bmatrix} 0 & A\Hh \\ A & 0 \end{bmatrix}.
\]



\section*{Problem 5. SVD application: orthogonal Procrustes}

Suppose the rows of $A,B\in\R^{m\times n}$ are corresponding observations in two coordinate
systems. The orthogonal Procrustes problem is
\[
  \min_{Q\T Q=I}\norm{AQ-B}_F. \tag{A.1}
\]

\begin{enumerate}[label=\textbf{A\arabic*.}]
\item Use the identity $\|M\|_F^2=\tr(M\T M)$ to expand the squared Frobenius norm
in (A.1), discard the $Q$-independent terms, and show that the problem is equivalent to
\[
   \max_{Q\T Q=I}\operatorname{tr}(Q\T A\T B). \tag{A.2}
\]
\item Let $A\T B=U\Sigma V\T$ be a full SVD. Prove that
\[
Q_\star=UV\T \tag{A.3}
\]
is a maximizer of (A.2). Your proof should introduce an orthogonal matrix $W$ and justify
the trace inequality used to bound $\operatorname{tr}(\Sigma W)$.
\item Make the equality case explicit. If $r=\rank(A\T B)$, show that all maximizers have the
form
\[
 Q=U\begin{bmatrix}I_r&0\\0&W_0\end{bmatrix}V\T, \tag{A.4}
\]
where $O(n-r)$ denotes the set of all $(n-r)\times(n-r)$ orthogonal matrices, i.e.,
$W_0\in\R^{(n-r)\times(n-r)}$ with $W_0\T W_0=I$. When $r=n$ the block $W_0$ is absent
($0\times 0$) and the formula reduces to $Q=UV\T$. Conclude precisely when the minimizer
is unique.
\item \textbf{Numerical experiment.} Choose $m=50$, $n=5$ (any $m\ge n$ works).
Generate a random $A\in\R^{m\times n}$ with i.i.d.\ $\mathcal{N}(0,1)$ entries,
an orthogonal $Q_0\in\R^{n\times n}$ (e.g., the $Q$ factor of the QR factorization
of a random Gaussian matrix), and a noise matrix $E\in\R^{m\times n}$ with i.i.d.\
$\mathcal{N}(0,1)$ entries. For each $\sigma\in\{0,10^{-3},10^{-2},10^{-1}\}$:
\begin{enumerate}[label=(\alph*)]
\item Set $B=AQ_0+\sigma E$.
\item Compute $Q_\star = UV\T$ from an SVD $A\T B = U\Sigma V\T$.
\item Report three Frobenius-norm residuals:
\[
\epsilon_{\mathrm{orth}}:=\norm{Q_\star\T Q_\star-I}_F,\qquad
\epsilon_{\mathrm{data}}:=\norm{AQ_\star-B}_F,\qquad
\epsilon_{\mathrm{align}}:=\norm{Q_\star-Q_0}_F.
\]
\end{enumerate}
\textbf{Note.} The $\sigma=0$ case is the noise-free sanity check $B=AQ_0$; there
$Q_\star$ should recover $Q_0$ to machine precision and both
$\epsilon_{\rm data},\epsilon_{\rm align}$ should be $\approx 0$.
\end{enumerate}

\section*{Problem 6. SVD application: canonical correlation analysis}

\textbf{Motivation.} Suppose $m$ samples carry measurements from two
blocks of variables: $X\in\R^{m\times p}$ (say, $p$ psychological test
scores) and $Y\in\R^{m\times q}$ (say, $q$ physical measurements), with
rows = samples, columns = variables (so the $j$-th column $X_j\in\R^m$
holds the $m$ observations of variable $j$), both blocks
\emph{column-centered} (i.e., $\frac{1}{m}\sum_{i=1}^m X_{ij}=0$ for
every column $j$). The ordinary Pearson correlation $\rho(X_j, Y_k)$
measures the linear association of a \emph{single} pair of variables,
but in practice variables within each block are correlated with one
another, and any block-level relationship is often split across many
variables rather than concentrated in one pair. \textbf{Canonical
correlation analysis (CCA)} asks for the strongest such
\emph{block-level} linear relationship: find linear combinations
$Xa, Yb$ (with $a\in\R^p$, $b\in\R^q$) that are maximally correlated.
Up to $\min(p, q)$ such pairs $(a_i, b_i)$ are produced, with
successive pairs constrained to be orthogonal to the previous ones so
that each new pair captures fresh block-level structure. The numbers
$\rho_i=\operatorname{Corr}(X a_i, Y b_i)$ are the \emph{canonical
correlations}, the vectors $a_i, b_i$ are \emph{canonical directions},
and the projections $X a_i, Y b_i$ are \emph{canonical variates}.

The Pearson correlation of two $m$-vectors $u, v$ (zero-mean, as
guaranteed by column-centering) is
$\operatorname{Corr}(u, v) = u\T v / (\norm{u}_2 \norm{v}_2)$, so the
single-pair CCA problem is
\[
 \max_{a\in\R^p,\, b\in\R^q}\, \operatorname{Corr}(Xa, Yb)
 \;=\; \max_{a, b}\, \frac{a\T (X\T Y) b}{\norm{Xa}_2\, \norm{Yb}_2}.
 \tag{B.1}
\]
Each remaining pair $(a_i, b_i)$, $i\ge 2$, maximizes (B.1) subject
to the additional constraints
$\operatorname{Corr}(X a_i, X a_j) = \operatorname{Corr}(Y b_i, Y b_j) = 0$
for every $j<i$.

Define the sample covariance matrices
\[
S_{xx}=\tfrac{1}{m-1}X\T X,\qquad
S_{yy}=\tfrac{1}{m-1}Y\T Y,\qquad
S_{xy}=\tfrac{1}{m-1}X\T Y.
\]
We assume throughout that $S_{xx}$ and $S_{yy}$ are positive definite,
and let $S_{xx}^{1/2}$ and $S_{yy}^{1/2}$ denote their symmetric
positive-definite square roots. Whitening by these square roots turns
the problem into one SVD.

\begin{enumerate}[label=\textbf{B\arabic*.}]
\item Derive a clean algebraic formulation of the single-pair problem
      (B.1) using $S_{xx}$, $S_{yy}$, $S_{xy}$.
\item Solve the single-pair formulation from B1. Give closed-form 
      expressions for the first canonical pair $(a_1, b_1)$ and the
      first canonical correlation $\rho_1$.
\item Generalize to all $k = \min(p, q)$ pairs. Give closed-form
      expressions for every $(a_i, b_i)$ and every $\rho_i$, and verify
      $a_i\T S_{xx} a_j = \delta_{ij}$ and
      $b_i\T S_{yy} b_j = \delta_{ij}$.
\end{enumerate}

\section*{Problem 7. Lecture 04 supplement}
\textit{See also: the Lecture~04 supplement at
\href{http://47.107.139.91:2601/files/lecture/lec04.pdf}{\texttt{http://47.107.139.91:2601/files/lecture/lec04.pdf}}.}

\subsection*{C.1 Max--min characterization (self-contained lemma)}
Let $A\in\R^{m\times n}$ with $m\ge n$ and singular values
$\sigma_1(A)\ge\cdots\ge\sigma_n(A)\ge0$. For $1\le i\le n$, define
\[
 \phi_i(A)=\max_{\substack{S\subseteq\R^n\\\dim S=i}}
       \;\min_{\substack{x\in S\\\norm{x}_2=1}}\norm{Ax}_2.
\]
\textbf{Task C1.} Prove that $\phi_i(A)=\sigma_i(A)$.

\subsection*{C.2 Perturbation bound for singular values}
\begin{theorem}[Weyl-type singular-value bound]
For $A,B\in\R^{m\times n}$, $m\ge n$,
\[
 \max_{1\le i\le n}\left|\sigma_i(A)-\sigma_i(B)\right|\le\norm{A-B}_2. \tag{C.1}
\]
\end{theorem}

\textbf{Task C2.} Prove Theorem 1. You may use the max--min characterization
$\phi_i(A)=\sigma_i(A)$ from Part C.1. As a sanity check, verify (C.1)
numerically on at least five random matrix pairs and report the largest
observed ratio $\max_i|\sigma_i(A)-\sigma_i(B)| / \norm{A-B}_2$.

\subsection*{C.3 Best rank-$k$ approximation and distance to rank deficiency}
Let $A\in\R^{n\times n}$ be nonsingular with SVD
$A=U\operatorname{diag}(\sigma_1,\ldots,\sigma_n)V\T$, where
$\sigma_1\ge\cdots\ge\sigma_n>0$.

\textbf{Task C3.} Fix $0\le k<n$. Determine the minimum-norm perturbation
that makes $A$ rank-deficient, taken over all $E\in\R^{n\times n}$:
\[
 \min_{\substack{\rank(A+E)\le k\\ E\in\R^{n\times n}}}\norm{E}_2. \tag{C.2}
\]

\end{document}
